\begin{abstract}
Detection benchmarks in aerial and remote-sensing settings are usually ranked by
mean average precision (mAP). We show that on a small, single-class landmine
surface-indicator dataset this ranking is uninformative for deployment: a
3.0\,M-parameter detector and an 11.1\,M-parameter detector differ by less than
the seed-to-seed spread in test mAP$_{50}$ ($0.979$ vs.\ $0.977 \pm 0.003$), yet their
recall-first operating points differ by a factor of $4.3$ in false positives per
image ($6.50$ vs.\ $1.51$). We argue that the meaningful object of study is the
operating point, not the summary curve, and we build the evaluation protocol
around that claim.

The study rests on two methodological contributions that require no new
architecture. First, we audit a widely reused landmine image corpus and find that
its published split leaks: $76$ Roboflow base-name groups span splits and
$63.8\%$ of evaluation images have a near-duplicate in training, consistent with
augmentation applied before splitting. We rebuild the split group-wise, apply a
two-method leak filter (perceptual hash plus pixel MSE) with manual review of
every candidate, and remove $30$ evaluation images. The frozen manifest
(3{,}907 / 181 / 189 images; digest \texttt{\manifestsum\ldots}) contains no
exact duplicate, no base-key collision, and no confirmed near-duplicate leak.

Second, we measure what the leak was worth. A detector trained on the leaky split
scores $0.994$ mAP$_{50}$ on the full clean test set against $0.974$ for a
clean-trained model, but on the $49$ contamination-free test images the two are
tied ($0.995$ vs.\ $0.995$). The apparent inflation is therefore contamination,
not capability---and it is small, a negative result we report rather than bury.

We then evaluate a MobileNetV3-Large second stage over detector crops on three
detector seeds. The two-stage arm improves precision over a recall-matched
single-stage baseline at two of three seeds ($\Delta P = +0.128$ and $+0.144$,
bootstrap intervals excluding zero) but fails at the third
($\Delta P = -0.037$, interval $[-0.077, +0.013]$). Under a max-$F_2$ rule a
val-tuned single-stage threshold beats the two-stage pipeline outright
($P = 0.976$ vs.\ $0.945$). We report the effect as detector-dependent rather
than architectural, and we identify the mechanism: the two-stage gain comes from
rescuing seeds whose val-selected threshold transferred poorly, not from better
localisation. Grad-CAM confirms that the detector's evidence lies on the objects
(median $58\times$ energy enrichment over chance) and that the second stage adds
almost no spatial evidence of its own ($1.82\times$ for the image-level baseline,
$2.13\times$ for crops)---a caution against reading a classifier's accuracy as
reasoning. All experiments were run on one 8\,GB laptop GPU within a
$\sim$30\,h budget; we state explicitly what that budget forbade.
\end{abstract}
